% Generated draft: requires mathematical and source review.
\documentclass[11pt]{article}
\usepackage[margin=25mm]{geometry}
\usepackage{fontspec}
\setmainfont{Latin Modern Roman}
\usepackage{amsmath,amssymb,mathtools,unicode-math}
\setmathfont{Latin Modern Math}
\usepackage{enumitem,needspace,xurl,hyperref}
\setlength{\parindent}{0pt}
\setlength{\parskip}{5pt}
\setlength{\emergencystretch}{3em}
\newcommand{\N}{\mathbb{N}}
\newcommand{\Z}{\mathbb{Z}}
\newcommand{\Q}{\mathbb{Q}}
\newcommand{\R}{\mathbb{R}}
\newcommand{\C}{\mathbb{C}}
\newcommand{\F}{\mathbb{F}}
\newcommand{\A}{\mathbb{A}}
\newcommand{\PP}{\mathbb{P}}
\newcommand{\E}{\mathbb{E}}
\newcommand{\eps}{\varepsilon}
\newcommand{\dd}{\,\mathrm{d}}
\begin{document}
\section*{UnsolvedMath ID: 30003973}
\textbf{Source problem code:} catalogue entry
\subsection{Definitions and mathematical statement}
The supplied record asks whether every pair of \(2\)-equivalent graphs is also \(3\)-equivalent, but does not define equivalence or specify the graph class. The following conditional formulation interprets equivalence as Ramsey equivalence, as suggested by the supplied reference title; these conventions require confirmation against the cited sources.

Under this interpretation, all graphs are finite, simple, and undirected. For graphs \(F,G\) and an integer \(r\geq 2\), write \(F\to_r G\) if every edge coloring \(c:E(F)\to\{1,\ldots,r\}\) contains a monochromatic copy of \(G\). A copy means a subgraph isomorphic to \(G\), not necessarily an induced subgraph; a coloring need not use every color. Define
\[
G\sim_r H
\quad\Longleftrightarrow\quad
\text{for every finite simple undirected graph }F,
\quad F\to_r G\ \Longleftrightarrow\ F\to_r H.
\]
The question, under these conventions, is whether
\[
\forall G\,\forall H,\qquad
G\sim_2 H\ \Longrightarrow\ G\sim_3 H.
\]
Here \(G\) and \(H\) range over all finite simple undirected graphs. The supplied record does not establish whether the intended formulation restricts this class, for example by excluding isolated vertices or requiring an edge. Consequently, this is a conditional statement-only draft, not a source-verified resolution of the missing definitions.
\subsection{Short English statement}
Under the stated Ramsey-equivalence interpretation, suppose exactly the same finite graphs force a monochromatic copy of \(G\) as force one of \(H\) in every two-color edge coloring. Must the same equivalence hold for three-color edge colorings? The intended graph class and equivalence definition require source confirmation.
\subsection{Sources}
\begin{itemize}
\item Supplied local registry record: permanent UnsolvedMath ID \(30003973\), source problem code \texttt{catalogue entry}, in the Oberwolfach Reports Open Problems collection. No individual proposer is supplied; the proposed year is recorded as 2018.
\item Supplied source citation: \emph{Mini-Workshop: Positional Games (2019)}, Oberwolfach Reports, DOI 10.4171/owr/2018/44, pp.\ 2715--2738. \url{https://doi.org/10.4171/owr/2018/44}. The supplied background describes Oberwolfach Reports as documenting mathematical workshops and their open-problem discussions.
\item Supplied literature reference: \emph{Ramsey equivalence}. \url{https://arxiv.org/abs/1809.09232}.
\item Literature assessment date: \texttt{2026-08-21}. The supplied dated curation labels the problem open and states that the implication from two equivalence to three equivalence remains posed as an open problem in the primary literature. It lists the preceding OWR report and Ramsey-equivalence reference as checked sources. This assessment is supplied dated curation, not independently verified here; the supplied background requests independent verification of time-sensitive status and cited claims.
\end{itemize}
\end{document}
